\section{Proposed Adaptation Methodology}

\subsection{Glottal Inverse Filtering (GIF) and Acoustic Normalization}
To address \textit{Melodic Tone Neutralization}, standard fundamental frequency ($F_0$) is insufficient because musical melodies overwrite pitch contours. We hypothesize that tonal contrasts in singing are physically preserved in the vocal tract excitation source. We apply Iterative Adaptive Inverse Filtering (IAIF) \cite{alku1992glottal,iseli2004f0} to perform source-filter deconvolution:
\begin{equation}
u_g'(t) = s(t) * h_{\text{vocal}}^{-1}(t) * h_{\text{lip}}^{-1}(t)
\end{equation}
where $s(t)$ is the acoustic signal, $h_{\text{vocal}}(t)$ represents time-varying vocal tract formants, and $h_{\text{lip}}(t)$ is the radiation impedance. The deconvolution applies pitch-synchronous analysis windows ($25$\,ms Hanning with $5$\,ms hop) to ensure robust vocal tract formant cancellation even when fundamental frequency $F_0 > 250$\,Hz. By integrating $u_g'(t)$ with a leaky factor of $0.99$, we recover the glottal volume velocity pulse $u_g(t)$ and compute four glottal phonation biomarkers:
\begin{enumerate}
    \item \textbf{Open Quotient ($O_q$):} The ratio of vocal fold opening duration to fundamental period $T_0$.
    \item \textbf{Formant-Corrected Spectral Tilt ($H_1^* - H_2^*$):} Computed using the Iseli-Alwan formulation \cite{iseli2004f0} to eliminate vowel formant amplification on harmonic amplitudes.
    \item \textbf{Normalized Amplitude Quotient (NAQ):} $\text{NAQ} = \frac{f_{\text{ac}}}{d_{\text{peak}} \cdot T_0}$ \cite{alku2002normalized}, characterizing glottal closing velocity.
    \item \textbf{Period-Doubling Ratio (PDR):} The subharmonic spectral energy ratio ($F_0 / 2$) detecting vocal fry and creaky glottalization.
\end{enumerate}

\subsection{Continuous F0 Conditioning and Syllable-Tone InfoNCE Loss}
To render phonetic representations invariant to pitch elevation while retaining lexical identity, we introduce a continuous pitch embedding module. Raw $F_0$ extracted via RMVPE \cite{zhou2023rmvpe} is mapped to semitone cents ($c = 1200 \log_2(F_0 / 10)$) and discretized into $256$ log-spaced bins. These are projected through a two-layer MLP to obtain pitch embedding $E_{F_0} \in \mathbb{R}^{T \times d_{\text{model}}}$, fused with acoustic encoder representations $Z_{\text{acoustic}}$ via learnable gating:
\begin{equation}
Z = Z_{\text{acoustic}} + \sigma(\gamma) \odot E_{F_0}
\end{equation}
where $\gamma$ is a trainable scalar parameter initialized to $0.0$.

To enforce melody invariance, we formulate a Syllable-Tone InfoNCE contrastive objective \cite{oord2018infonce}:
\begin{equation}
\mathcal{L}_{\text{InfoNCE}} = -\log \frac{\exp(\text{sim}(z_i^{\text{sing}}, z_i^{\text{speech}}) / \tau)}{\sum_{j=1}^K \exp(\text{sim}(z_i^{\text{sing}}, z_j^{\text{speech}}) / \tau)}
\end{equation}
where $\text{sim}(u, v) = \frac{u^\top v}{\|u\| \|v\|}$, and temperature $\tau = 0.07$. This pulls representations of identical syllables sung and spoken closer together in latent space while pushing dissimilar syllables apart.

\subsection{Multi-Dialect LoRA Mixture-of-Experts (MoE) Architecture}
Vietnamese singing features pronounced phonological divergence across three major regional dialects: Northern (Hanoi: rigid 6-tone distinction), Central (Hue: shortened vowels, heavy glottal drop), and Southern (Saigon: \textit{H\H{o}i}--\textit{Ng\~a} merger, final consonant shifts $-n/-ng \to -\eta$). 

We implement a sparsely-gated Mixture-of-Experts architecture \cite{shazeer2017outrageously,fedus2022switch} combined with Low-Rank Adaptation (LoRA) \cite{hu2021lora}. Three regional expert adapters ($\text{LoRA}_{\text{N}}$, $\text{LoRA}_{\text{C}}$, and $\text{LoRA}_{\text{S}}$ with rank $r=8, \alpha=16$) are integrated into the multi-head attention projections. An Acoustic Dialect Router computes softmax routing probabilities:
\begin{equation}
\begin{aligned}
g(x) &= \text{Softmax}(W_g \cdot \bar{z}_{\text{enc}} + b_g) \\
\Delta W &= \sum_{k=1}^3 g_k(x) \cdot \left(\frac{\alpha}{r} B_k A_k\right)
\end{aligned}
\end{equation}
To prevent expert collapse, we optimize the expert load-balancing loss:
\begin{equation}
\mathcal{L}_{\text{balance}} = \alpha_{\text{bal}} \cdot N \sum_{k=1}^N m_k P_k
\end{equation}
where $m_k$ denotes the fraction of tokens routed to expert $k$, and $P_k$ is the average routing probability. Crucially, while colloquial Southern speech merges tones \textit{H\H{o}i} and \textit{Ng\~a} into a unified dipping contour, Southern vocalists frequently sing standardized literary lyrics with canonical orthography. The Southern expert $\text{LoRA}_{\text{S}}$ is trained under auxiliary cross-entropy to project colloquial merged acoustics directly onto standard orthographic targets, eliminating false homophone confusions without requiring separate post-processing disambiguation.

\subsection{Compound Objective and Random Search Champion \texorpdfstring{($RS_{01}$)}{(RS01)}}
The complete network is optimized under a 6-component multi-task objective:
\begin{equation}
\begin{aligned}
\mathcal{L}_{\text{total}} &= w_{\text{ce}} \mathcal{L}_{\text{ce}} + w_{\text{diag}} \mathcal{L}_{\text{diag}} + w_{\text{mono}} \mathcal{L}_{\text{mono}} \\
&\quad + w_{\text{consist}} \mathcal{L}_{\text{consist}} + w_{\text{entropy}} \mathcal{L}_{\text{entropy}} + w_{\text{replay}} \mathcal{L}_{\text{replay}}
\end{aligned}
\end{equation}
Here, $\mathcal{L}_{\text{consist}}$ enforces Jensen-Shannon divergence consistency between polyphonic mix predictions and clean vocal stem predictions; $\mathcal{L}_{\text{entropy}}$ penalizes cross-attention blur over sustained vowels; and $\mathcal{L}_{\text{replay}}$ preserves spoken conversational anchors. Through multi-objective random search over the Pareto frontier:
\begin{equation}
\begin{aligned}
\mathcal{S}_{\text{Pareto}} &= 0.40 \cdot \text{WER}_{\text{poly}} + 0.30 \cdot \text{WER}_{\text{vocal}} \\
&\quad + 0.30 \cdot \text{WER}_{\text{speech}}
\end{aligned}
\end{equation}
our champion configuration ($RS_{01}$) establishes optimal weights: $w_{\text{ce}} = 0.943$, $w_{\text{diag}} = 1.817$, $w_{\text{mono}} = 2.189$, $w_{\text{consist}} = 1.455$, $w_{\text{entropy}} = 0.042$, and $w_{\text{replay}} = 1.176$.

\subsection{Low-Latency Streaming Monotonic Chunkwise Attention (MoChA)}
Standard full-sequence attention incurs $O(T^2)$ lookahead complexity, rendering it unusable for live karaoke or real-time subtitling. We integrate Monotonic Chunkwise Attention (MoChA) \cite{chiu2018monotonic}. At each step, a monotonic energy layer computes Bernoulli stepping probabilities:
\begin{equation}
e_{i, j} = v^\top \tanh(W_q s_{i-1} + W_k h_j + b), \quad p_{i, j} = \sigma(e_{i, j})
\end{equation}
Online decoding advances monotonically to frame $j$ upon $p_{i, j} > 0.45$. Softmax chunk attention is then restricted to a localized window of $C = 160$\,ms ($8$ frames) with $40$\,ms lookahead:
\begin{equation}
\begin{aligned}
c_i &= \sum_{k=j-C+1}^j \beta_{i, k} h_k \\
\beta_{i, k} &= \frac{\exp(q_i^\top k_k / \sqrt{d})}{\sum_{m=j-C+1}^j \exp(q_i^\top k_m / \sqrt{d})}
\end{aligned}
\end{equation}

\begin{figure}[t]
\centering
\includegraphics[width=0.92\linewidth]{cross_attention_alignment_comparison.png}
\caption{Cross-attention alignment maps: standard Whisper attention matrix stalls horizontally over prolonged vibrato vowels (left), whereas monotonic chunkwise attention enforces strictly sequential temporal progression without looping (right).}
\label{fig:cross_attention}
\end{figure}

As illustrated in Fig.~\ref{fig:cross_attention}, this monotonic formulation bounds computation latency to $8.50$\,ms per chunk. Adding the $40$\,ms lookahead window, the combined algorithmic latency is $48.50$\,ms, strictly below the $100$\,ms human perception threshold for real-time interaction.

\subsection{Fast Syllable Adaptation for Rapid Rhythmic Delivery}
In modern musical genres such as Vietnamese rap and fast melodic flows, speech rate frequently exceeds $8$ syllables per second (SPS), causing standard frame downsamplers to blur brief consonant onsets. We introduce the \textit{Fast Syllable Adapter} (FSA, $1.219$\,M parameters), which incorporates transient-preserving multi-resolution temporal pooling:
\begin{equation}
\begin{aligned}
Z_{\text{fast}} &= \text{LayerNorm}\Bigl(Z + \sum_{k \in \{1, 2, 4\}} \text{Conv1D}_k\bigl(\text{GELU}( \\
&\quad \text{DepthwiseConv}_k(Z))\bigr)\Bigr)
\end{aligned}
\end{equation}
By executing parallel temporal convolutions with kernel sizes $k \in \{1, 2, 4\}$ and stride $1$, the adapter preserves sub-frame consonant transients while capturing rhythmically compressed vowel contours.

\subsection{Multi-Speaker Duet Disentanglement (D-SADA)}
When lead and backing vocals overlap simultaneously, single-stream decoders suffer from severe amplitude dominance, deleting the secondary vocal. We propose the \textit{Dual-Stream Attentive Disentanglement Adapter} (D-SADA, $5.804$\,M parameters). Utilizing multi-pitch harmonic gates conditioned on dual $F_0$ tracks ($F_{0,1}, F_{0,2}$), the module projects acoustic features into orthogonal lead and secondary streams ($Z_{\text{lead}}, Z_{\text{sec}}$) optimized via utterance-level Permutation Invariant Training (uPIT) \cite{kolbaek2017pit}:
\begin{equation}
\mathcal{L}_{\text{uPIT}} = \min_{\pi \in \mathcal{P}} \sum_{s=1}^2 \mathcal{L}_{\text{CE}}(\hat{y}_{\pi(s)}, y_s)
\end{equation}
where $\mathcal{P}$ represents the set of all permutations over the two vocal tracks.

\subsection{Instrument-Aware Tri-Band Suppression Architecture (IAS-Adapter)}
Polyphonic musical accompaniments collide with singing formants across distinct frequency bands depending on instrument physics. We design the \textit{Instrument-Aware Suppression Adapter} (IAS-Adapter, $3.236$\,M parameters), which decomposes latent acoustic representations into three functional sub-bands:
\begin{enumerate}
    \item \textbf{Low-Band Gate ($<250$\,Hz):} Conditioned on 808/synth bass detection probability $P_{\text{synth}}$, attenuating sub-bass floor leakage.
    \item \textbf{Mid-Band Formant Gate ($250\text{--}3{,}000$\,Hz):} Self-attention cross-projection conditioned on $[P_{\text{piano}}, P_{\text{guitar}}]$ that desmears $F_1, F_2$ vowel formants against polyphonic chords and vibrato modulations.
    \item \textbf{High-Band Transient Limiter ($>3{,}000$\,Hz):} Gated soft limiter activated by drum probability $P_{\text{drum}}$, dampening snare/hi-hat transients while preserving voiceless consonant edges (/t/, /th/, /s/, /x/).
\end{enumerate}
The suppressed sub-bands are recombined via projection and a residual identity bypass:
\begin{equation}
\label{eq:ias_recon}
Y = \text{LayerNorm}(X + \gamma_{\text{res}} \cdot \text{Proj}([Z_{\text{low}}', Z_{\text{mid}}', Z_{\text{high}}']))
\end{equation}
where $\gamma_{\text{res}}$ is initialized to $0.20$ to maintain baseline representation integrity.

\subsection{Edge INT8 Post-Training Quantization}
For deployment on resource-constrained consumer edge devices, we implement dynamic-range post-training INT8 quantization \cite{jacob2018quantization}. Linear feed-forward projections and attention projections within our adapter suite are quantized dynamically from 32-bit floating point ($\text{FP32}$) to signed 8-bit integers ($\text{qint8}$):
\begin{equation}
q = \text{clamp}\left(\left\lfloor \frac{w}{S} \right\rceil + Z_{\text{zero}}, -128, 127\right)
\end{equation}
where $S$ is the dynamic scale factor and $Z_{\text{zero}}$ is the zero-point offset. LayerNorm, Softmax, and residual additions are maintained in FP32 to ensure numerical stability and prevent gradient degradation.

\subsection{Unified End-to-End Pipeline and Dynamic Scenario Gating}
In practical production deployments, the four acoustic adapter modules operate as a hierarchical, scenario-aware cascaded pipeline rather than isolated components. Raw input audio first enters the \textbf{IAS-Adapter} ($Z_{\text{low}}, Z_{\text{mid}}, Z_{\text{high}}$) to decouple transient percussion, piano formant smearing, and sub-bass energy. If the input is detected as clean vocal or spoken speech, the residual gate ($\gamma_{\text{res}}$ in Eq.~\ref{eq:ias_recon}) automatically attenuates sub-band suppression, passing acoustic frames unaltered. Next, if multi-pitch tracking indicates simultaneous harmonic voices, \textbf{D-SADA} activates dual-stream uPIT projection; otherwise, single-lead features bypass directly. The representation is then conditioned by the \textbf{Fast Syllable Adapter (FSA)} whenever rhythmic pacing exceeds $5.0$\,SPS. Finally, continuous pitch embedding ($E_{F_0}$) and the Sparse Dialect MoE router ($g_k(x)$) modulate the encoder representations before the low-latency \textbf{MoChA} streaming decoder decodes tokens frame-by-frame. As quantified in Table~\ref{tab:edge_quantization}, the complete integrated pipeline consumes only $21.82$\,MB in INT8 and executes in $64.31$\,ms ($1.31\times$ real-time latency on CPU), achieving an INT8 Real-Time Factor of $0.02108$ ($47.4\times$ faster than real-time).

